در شکل (c) با \lr{zero-padding} عملاً تصویر اصلی داخل یک ناحیهٔ بزرگ‌تر قرار گرفته و بیرون آن صفر شده است؛ یعنی
$$
g(x,y)=f(x,y)\,w(x,y),
\qquad
w(x,y)=\mathrm{rect}\!\Big(\frac{x}{L_x}\Big)\mathrm{rect}\!\Big(\frac{y}{L_y}\Big)
$$
که $w$ یک پنجرهٔ مستطیلیِ هم‌راستا با محورهای افقی/عمودی است. در حوزهٔ فوریه داریم
$$
G(u,v)=F(u,v)\ast W(u,v),
\qquad
W(u,v)=L_xL_y\,\mathrm{sinc}(L_xu)\,\mathrm{sinc}(L_yv).
$$
پس طیفِ تصویرِ جدید برابر کانولوشنِ طیفِ تصویرِ اصلی با طیفِ پنجرهٔ مستطیلی است. چون $W(u,v)$ جداپذیر و دارای لوب‌های قوی در راستای محورهای $u$ و $v$ است، این کانولوشن باعث می‌شود انرژیِ طیف در امتداد $u=0$ و $v=0$ (محورهای عمودی و افقیِ شکل (d)) تقویت و به صورت یک «صلیب» روشن دیده شود.

تعبیر فضاییِ همین پدیده این است که صفر کردنِ بیرونِ تصویر، در مرزهای ناحیهٔ تصویر چهار لبهٔ تیزِ افقی/عمودی (تغییر ناگهانی شدت) ایجاد می‌کند. لبهٔ عمودی تقریباً در راستای $y$ ثابت است، لذا مؤلفه‌های فرکانسی‌اش عمدتاً روی محور $v=0$ متمرکز می‌شوند و لبهٔ افقی نیز انرژی را روی محور $u=0$ زیاد می‌کند؛ بنابراین نسبت به شکل (b) افزایش قابل توجه روی محورهای افقی و عمودیِ شکل (d) مشاهده می‌شود.